\vspace{0.25cm}
La partie « Systèmes ouverts en régime stationnaire » complète la partie « Machines thermique» du programme de première année de la classe de PCSI en proposant notamment un bilan d'entropie. La maîtrise des démonstrations par les étudiants et l'application des résultats à des situations concrètes constituent des objectifs de formation.

\begin{tabularx}{\textwidth}[t]{|X|X|}
\hline \textbf{Notions et contenus} & \textbf{Capacités exigibles} \\
\hline 3.1. Systèmes ouverts en régime stationnaire \\
\hline Premier et deuxième principes de la & Établir les relations $\Delta h+\Delta e=w_u+q$ et \\
thermodynamique pour un système ouvert en & $\Delta s=s_e+s_c$ et les utiliser pour étudier des \\
régime stationnaire, dans le seul cas d'un & machines thermiques réelles à l'aide de \\
écoulement unidimensionnel au niveau de la & diagrammes thermodynamiques ($T$, $s$) et \\
section d'entrée et de la section de sortie. & $(p, h)$. \\
\hline
\end{tabularx}
